\section{2D isosceles quantum Schr\"odinger equation}
\label{sec:2d_isosceles}

\subsection{Motivation}

Section~\ref{sec:triangle_stability} (ex30) showed that the equilateral
triangle is a \textbf{saddle point} in shape space: the antisymmetric
deformation mode has a negative Hessian eigenvalue
$\lambda_{\rm sh} = -0.117\,E_{\rm Pl}$.
This means the 1D equilateral breathing-mode calculation of ex23--ex29
restricts to a saddle of the potential, giving \emph{upper bounds} on $E_0$.
The question is: by how much does the true ground-state energy differ?

To answer this, we implement the full 2D quantum Schr\"odinger equation
on the \emph{isosceles submanifold} $d_{12}=d_{13}=a$, $d_{23}=b$
(the natural complement to the equilateral line), and compare with the 1D result.

\subsection{Classical isosceles minimum (ex32)}

Before solving the quantum problem, we locate the classical energy minimum
on the isosceles submanifold using the effective potential
$E_{\rm eff}(a,b) = \text{ZP}(\bar{d}) + V_2(a,a,b) + V_3(a,a,b)$
where $\bar{d}=(2a+b)/3$.

\begin{table}[h]
\centering
\caption{Classical isosceles minimum vs equilateral minimum at $C=C_{\rm Pl}$,
$m_{\rm sp}=0.1\,l_{\rm Pl}^{-1}$.  The isosceles minimum lies at $b\to 0$
(Yukawa singularity), confirming that the effective-potential landscape has no
classical lower bound off the equilateral line.}
\label{tab:iso_classical}
\begin{tabular}{lrrrr}
\hline
Configuration & $a$ & $b$ & $b/a$ & $E$ [$E_{\rm Pl}$] \\
\hline
Equilateral min.\ & $1.371$ & $1.371$ & $1.000$ & $-0.265$ \\
Isosceles (grid) & $2.293$ & $0.300$ & $0.131$ & $-0.446$ \\
True limit $b\to0$ & any & $0$ & $0$ & $-\infty$ \\
\hline
\end{tabular}
\end{table}

\paragraph{Physical interpretation.}
The Yukawa pair potential $V_2 = -C^2 e^{-m_{\rm sp}b}/b \to -\infty$ as $b\to0$.
The effective ZP term $3/(8\bar{d}^2)$ is finite as $b\to0$ (with $a$ fixed),
so the effective potential is \emph{unbounded below} in the isosceles direction.

This is a classical collapse singularity: the 3-body Yukawa potential has no
hard core. The equilateral ``minimum'' found in ex23--ex29 is a local minimum
only within the equilateral subspace $d_{12}=d_{13}=d_{23}$.

Quantum mechanics resolves the collapse: the kinetic energy operator
$\sim -\nabla^2$ diverges as $b\to0$ faster than the Yukawa attraction
$\sim -1/b$, providing a quantum-mechanical hard core.
The correct treatment therefore requires the full quantum kinetic energy,
not the semi-classical ZP approximation.

\subsection{2D Jacobi coordinates (ex33)}

\paragraph{Coordinate system.}
For 3 particles of mass $m=m_{\rm Pl}$ in 1D, the Jacobi coordinates are:
\begin{align}
  \boldsymbol{\xi}_1 &= \mathbf{r}_2 - \mathbf{r}_3, \quad
    |\boldsymbol{\xi}_1| = b = d_{23},\\
  \boldsymbol{\xi}_2 &= \mathbf{r}_1 - \tfrac{1}{2}(\mathbf{r}_2+\mathbf{r}_3),\quad
    |\boldsymbol{\xi}_2| = h = \sqrt{a^2 - b^2/4},
\end{align}
with reduced masses $\mu_1 = m/2 = 0.5$ and $\mu_2 = 2m/3 \approx 0.667$.

\paragraph{Hamiltonian.}
The quantum Hamiltonian in $(b,h)$ coordinates (no ZP term — kinetic energy
is exact):
\begin{equation}
  \hat{H} = -\frac{1}{2\mu_1}\frac{\partial^2}{\partial b^2}
            -\frac{1}{2\mu_2}\frac{\partial^2}{\partial h^2}
            + V(b,h)
  = -\frac{\partial^2}{\partial b^2}
    -\frac{3}{4}\frac{\partial^2}{\partial h^2}
    + V(b,h),
\end{equation}
where the isosceles potential is
\begin{equation}
  V(b,h) = V_2\!\left(a,a,b\right) + V_3\!\left(a,a,b\right),
  \quad a = \sqrt{h^2 + b^2/4}.
\end{equation}
No separate ZP term appears: the full quantum kinetic energy provides
the natural repulsion against collapse.

\paragraph{Kinetic energy for the equilateral breathing mode.}
Along the equilateral line $b=d$, $h=d\sqrt{3}/2$, the effective kinetic
energy coefficient for the breathing coordinate $d$ is
\begin{equation}
  T_{\rm br} = -\left[\frac{1}{2\mu_1}\left(\frac{db}{dd}\right)^2
                + \frac{1}{2\mu_2}\left(\frac{dh}{dd}\right)^2\right]
               \frac{d^2}{dd^2}
  = -\left[1 + \frac{3}{4}\cdot\frac{3}{4}\right]\frac{d^2}{dd^2}
  = -\frac{25}{16}\frac{d^2}{dd^2},
\end{equation}
corresponding to effective mass $\mu_{\rm eff} = 8/25 = 0.32\,m_{\rm Pl}$.
The 1D equilateral calculation (ex23--ex29) used $\mu=1\,m_{\rm Pl}$,
which is $3.1\times$ too large — it \emph{underestimates} the kinetic energy,
overestimating binding.

\subsection{Numerical results}

\paragraph{Grid.}
We use a uniform $(b,h)$ grid with $N_b\times N_h = 50\times50$ points,
$b,h\in[0.15,\,20]\,l_{\rm Pl}$, and Dirichlet boundary conditions
($\psi=0$ on all edges).  The sparse Hamiltonian
($2500\times2500$, diagonal + nearest-neighbour off-diagonals) is diagonalised
with \texttt{scipy.sparse.linalg.eigsh}.

\begin{table}[h]
\centering
\caption{Ground-state energy comparison: 1D equilateral vs 2D Jacobi isosceles.
At $C=C_{\rm Pl}$, $m_{\rm sp}=0.1$.}
\label{tab:2d_iso_comparison}
\begin{tabular}{lrrrr}
\hline
Method & $E_0$ [$E_{\rm Pl}$] & $\langle b\rangle$ & $\langle a\rangle$ & $b/a$ (peak) \\
\hline
1D equilateral (ex23--ex29) & $-0.0698$ & --- & $4.52$ & $1.000$ \\
2D Jacobi isosceles (ex33) & $-0.0088$ & $6.88$ & $7.53$ & $0.956$ \\
\hline
Difference $\Delta E$ & $+0.0610$ ($87\%$ of $|E_0^{\rm 1D}|$) & & & \\
\hline
\end{tabular}
\end{table}

\begin{table}[h]
\centering
\caption{Ground-state energy vs $m_{\rm sp}$ from 2D Jacobi calculation
($N=40\times40$, $C=C_{\rm Pl}$).  For $m_{\rm sp}\geq0.130$, the 2D
calculation gives an unbound state ($E_0>0$).}
\label{tab:2d_iso_scan}
\begin{tabular}{rrrrr}
\hline
$m_{\rm sp}$ [$l_{\rm Pl}^{-1}$] & $\lambda$ [$l_{\rm Pl}$] &
  $E_0^{\rm 1D}$ & $E_0^{\rm 2D}$ & $\Delta E$ [$E_{\rm Pl}$] \\
\hline
$0.076$ & $13.2$ & $-0.117$ & $-0.049$ & $+0.068$ \\
$0.100$ & $10.0$ & $-0.070$ & $-0.013$ & $+0.057$ \\
$0.130$ & $\ 7.7$ & $-0.034$ & $+0.010$ & $+0.044$ \\
$0.150$ & $\ 6.7$ & $-0.018$ & $+0.019$ & $+0.037$ \\
$0.180$ & $\ 5.6$ & $-0.005$ & $+0.026$ & $+0.031$ \\
$0.198$ & $\ 5.1$ & $\approx0$ & $+0.029$ & $+0.029$ \\
\hline
\end{tabular}
\end{table}

\subsection{Physical interpretation}

\paragraph{Wavefunction geometry.}
At $m_{\rm sp}=0.1$, $C=C_{\rm Pl}$, the ground-state wavefunction peaks at
$(b,h)=(5.01,\,4.61)\,l_{\rm Pl}$, corresponding to
$a=\sqrt{4.61^2+5.01^2/4}=5.24\,l_{\rm Pl}$ and
$b/a = 5.01/5.24 = 0.956$.  The equilateral value is $b/a=1$, so the peak
is only $4.4\%$ off the equilateral line.  The quantum ground state is
effectively equilateral in shape.

\paragraph{Kinetic energy discrepancy.}
The 2D Jacobi calculation gives $E_0^{\rm 2D} = -0.0088$ vs
$E_0^{\rm 1D}=-0.0698$.  The discrepancy ($\Delta E = +0.061$, 87\% of
$|E_0^{\rm 1D}|$) arises primarily from the effective mass:
\begin{itemize}
  \item 1D equilateral (ex23): effective mass $\mu_{\rm eff}=1\,m_{\rm Pl}$
    (kinetic coefficient $1/2$), ZP $= 3/(8d^2)$.
  \item 2D Jacobi (ex33): kinetic coefficients $1$ (b-direction)
    and $3/4$ (h-direction), no ZP term.
  \item Along equilateral breathing mode: the 2D kinetic coefficient is
    $25/16 = 1.5625$, versus $1/2$ in the 1D model — a factor $3.1\times$ larger.
\end{itemize}
The larger kinetic energy in the 2D Jacobi model provides stronger quantum
pressure, reducing binding by a factor $\sim 3$.

\paragraph{Efimov window in 2D Jacobi.}
The 2D Jacobi calculation predicts the 3-body bound state threshold at
$m_{\rm sp}\approx0.12\,l_{\rm Pl}^{-1}$ (below which $E_0<0$, above which
$E_0>0$), compared with $m_{\rm sp}=0.198$ in the 1D equilateral model.
This shift reflects the stronger kinetic energy in the correct Jacobi treatment.

\subsection{Conclusion and open problems}

The 2D Jacobi Schr\"odinger equation for the isosceles submanifold gives
a less-bound ground state than the 1D equilateral approximation, due to the
correct Jacobi reduced masses.  Key conclusions:
\begin{enumerate}
  \item The ground-state wavefunction is near-equilateral in shape ($b/a\approx0.96$),
    confirming the equilateral as the preferred geometry.
  \item The equilateral approximation of ex23--ex29 \emph{overestimates} binding
    because it uses an effective mass $\mu=1$ instead of the correct
    $\mu_{\rm eff}=0.32$ for the equilateral breathing mode.
  \item The corrected (2D Jacobi) ground-state energy is
    $E_0^{\rm 2D}=-0.0088\,E_{\rm Pl}$ at $C_{\rm Pl}$, $m_{\rm sp}=0.1$.
  \item The Efimov window upper boundary shifts from $m_{\rm sp}=0.198$ (1D)
    to $m_{\rm sp}\approx0.12$ (2D Jacobi).
  \item A full 3D calculation including all angular degrees of freedom (not just
    the isosceles submanifold) is required for the physically correct Efimov
    spectrum.  This is the primary open computational problem (Q5).
\end{enumerate}

\paragraph{Status of Q5 (full 2D quantum Efimov).}
The 2D Jacobi calculation of ex33 gives a first quantitative estimate of the
true (beyond equilateral-approximation) quantum ground state.  The result:
$E_0^{\rm true}\approx -0.009\,E_{\rm Pl}$ at $C_{\rm Pl}$, $m_{\rm sp}=0.1$,
and the Efimov window narrows to $m_{\rm sp}\lesssim0.12\,l_{\rm Pl}^{-1}$.
The next step is a full 3D Faddeev or hyperspherical calculation.
\label{sec:q5_status}
